\section{Introduction}

Agent evaluation is quickly moving beyond single-turn QA into web navigation,
mobile environments, software engineering, and repository-grounded tool use
\cite{agentbench,webarena,visualwebarena,androidworld,swebench,swelancer,ambigswe,longsoftwaretasks,gittaskbench,toolgym,deepresearchbench,mcpbench,mcpuniverse}.
Real deployments are rarely pure language tasks. They require interacting with
interfaces, repositories, and heterogeneous tools
while producing structured outputs that can be checked automatically.

At the same time, the dominant reporting habit has remained simple. A model is
usually assigned one benchmark score, often one success rate, and the result is
read as an intrinsic capability measure. Recent work has already pushed back on
that simplification from several angles. Trajectory-oriented analysis argues
that final success hides diagnostically important execution structure
\cite{agentrx,trajeval}. Measurement work on long-context and benchmark scaling
argues that aggregate scores can conceal important task-specific variation
\cite{lostmiddle,longbench,ruler,infinitebench,nolima,loogle,illusiondiminishingreturns}.
Multimodal evaluation further shows that answer channel, extraction quality, and
benchmark family can all matter for what a score actually means
\cite{mme,mmmu,mathvista,hallusionbench,blink}. These critiques point in the
same direction: one reported score can hide multiple latent conditions of
success.

This paper studies a narrow measurement problem in \gittaskbench.
\gittaskbench already separates \processmetric from \resultmetric. The former
marks traces that reach grading with an output artifact. The latter marks
traces whose recorded artifact satisfies the official task criterion. Some
traces never reach grading; others reach grading but still fail. If both cases
are collapsed into a single final success rate, important structure
disappears.

The empirical design stays deliberately simple. The primary evidence is a
comparable open-model cohort on \gittaskbench. We then add a smaller
\texttt{Trafilatura\_01/02} subset on web-extraction tasks that extends
coverage to provider models, including \texttt{gpt-5.2} and
\texttt{Gemini-2.5-Pro}. We also add a targeted provider cohort on
\texttt{gpt-5-chat}, \texttt{gpt-5.2}, and \texttt{Gemini-2.5-Pro} so the
paper does not reduce to a purely within-open-model comparison. For the
task-aligned \texttt{Trafilatura} subset, we align identical benchmark tasks
across models without treating that subset as a frontier ranking. All plots and
tables report official \processmetric and \resultmetric totals or deterministic
aggregates derived from them. Finally, we add full-benchmark provider sweeps for
\texttt{gpt-5-chat}, \texttt{gpt-5.2}, and \texttt{Gemini-2.5-Pro} to test
whether the targeted subset remains representative once the full benchmark
distribution is restored.

The evidence hierarchy is intentional. The comparable cohort is the paper's
primary evidence. The \texttt{Trafilatura} subset, the targeted provider
cohort, the supplementary full-benchmark provider cohorts, and the short
open-model breadth extension are secondary analyses used to test whether the
same family structure survives broader model coverage.

Three findings motivate the paper. First, official success is narrow and
domain-skewed. In the comparable cohort, all observed \resultmetric mass
lies in office, security, web, and speech families; image and video families
are complete zeros. Second, tied headline totals hide family specialization.
The three headline leaders in the comparable cohort tie on official
\resultmetric success, but they succeed
on different tasks and families. The targeted provider cohort adds a second
measurement point: broader frontier coverage changes which families activate,
but provider models are themselves highly non-uniform on the same official
tasks. The full-benchmark provider sweeps sharpen that point further:
\texttt{gpt-5.2} expands into additional speech, security, physiological-signal,
and image families, \texttt{gpt-5-chat} remains concentrated in office, web,
and security with only a narrow image exception, and
\texttt{Gemini-2.5-Pro} stays narrow. Third,
\processmetric$=$True,
\resultmetric$=$False rows are not random residue. They cluster in structured
extraction families and mix threshold misses, parser failures, and
evaluation-side errors after grading begins.

The supported claim is simple. On the evaluated \gittaskbench cohorts,
reporting only a single global \resultmetric rate hides family-specific
competence and a structured \processmetric$=$True,
\resultmetric$=$False regime already visible in the benchmark's own outputs.

\paragraph{Contributions.}
This paper makes four concrete contributions.
\begin{enumerate}[leftmargin=1.25em]
\item We construct an official-metric study of \gittaskbench over a
comparable open-model cohort together with task-aligned and provider-side
supplementary cohorts, while keeping the benchmark's official
\processmetric/\resultmetric contract fixed.
\item We show that official \resultmetric success is highly concentrated in a
small subset of domains and families rather than being smoothly distributed
across the benchmark.
\item We show that tied global \resultmetric totals hide material
task-family specialization across models.
\item We identify an official \processmetric$=$True, \resultmetric$=$False
regime that clusters in structured extraction families and should be reported
explicitly in future repo-grounded evaluations.
\end{enumerate}

The paper is therefore a descriptive benchmark study. We report exact observed
\processmetric/\resultmetric counts from the evaluated cohorts, and we use the
provider slices as supplementary analyses rather than as claims about a stable
frontier ranking.
